% GENERATED FILE. Edit canonical lessons, then run npm run generate:books.
% canonical-source-hash: fnv1a64:fd665d7fcb314978
% canonical-lessons: UR-C411-dhalan, UR-C411-lurhakna, UR-C411-baramada, UR-C411-tah-khana, UR-C411-baghal

\chapter{Slopes, Verandas and Cellars}
\label{ch:ur-slopes-verandas-and-cellars}
\hlchaptermodality{411}

\begin{chapteropening}
\textbf{By the end of this chapter:} \emph{I can describe the lie of a place: what slopes, what is beside it, what is below.}

Complete the last lesson of chapter 411: I can describe the lie of a place: what slopes, what is beside it, what is below.
\end{chapteropening}

\section[\texorpdfstring{\ur{ڈھلان} (ḍhalān)}{ڈھلان (ḍhalān)}]{\ur{ڈھلان} (ḍhalān) — a slope}

\label{lesson:UR-C411-dhalan}



\begin{quote}
Before the new one: say the Urdu for a pole, then the Urdu for to show.
\end{quote}

\subsection*{You'll want to know: \ur{ڈھلان}}
\textbf{\ur{ڈھلان}} — \emph{ḍhalān} — \textquotedblleft{}a slope\textquotedblright{}.

\begin{quote}
\ur{پہاڑی} \ur{کی} \ur{ڈھلان} \ur{پر} \ur{گھاس} \ur{ہے۔}
\end{quote}

\begin{quote}
\emph{pahāṛī kī ḍhalān par ghās hai.}
\end{quote}

\begin{quote}
There is grass on the slope of the hill.
\end{quote}

Feminine.

\begin{cousinweb}
\textbf{\ur{ڈھلان}} (\emph{ḍhalān}) is built on \emph{ḍhalnā}, \textquotedblleft{}to topple, to slope away\textquotedblright{}, which CDIAL lists under \emph{ḍhalati}, \textquotedblleft{}bends over, falls\textquotedblright{}, a form it reconstructs (CDIAL 5581).
\end{cousinweb}

\begin{grammarlens}[title={What you've built}]
One more word for getting around.
\end{grammarlens}

\subsection*{Your turn}
\begin{itemize}
\raggedright
  \item \emph{Say it:} \emph{ḍhalān}
  \item \emph{Say it:} \emph{ḍhalān}, once more
  \item \emph{Say it:} \emph{pahāṛī kī ḍhalān par ghās hai.} — There is grass on the slope of the hill.
  \item \emph{Recall:} say the Urdu for a pole, then the Urdu for to show, then say \emph{ḍhalān} again
\end{itemize}

\begin{tcolorbox}[breakable,colback=teal!4,colframe=teal!35!black,title={Before you move on}]
What does \ur{ڈھلان} mean? (\textquotedblleft{}A slope\textquotedblright{}.) Say it once more.
\end{tcolorbox}
\section[\texorpdfstring{\ur{لڑھکنا} (luṛhaknā)}{لڑھکنا (luṛhaknā)}]{\ur{لڑھکنا} (luṛhaknā) — to roll (down)}

\label{lesson:UR-C411-lurhakna}



\begin{quote}
Before the new one: say the Urdu for to show, then the Urdu for a slope.
\end{quote}

\subsection*{You'll want to know: \ur{لڑھکنا}}
\textbf{\ur{لڑھکنا}} — \emph{luṛhaknā} — \textquotedblleft{}to roll (down)\textquotedblright{}.

\begin{quote}
\ur{ایک} \ur{پتھر} \ur{ڈھلان} \ur{سے} \ur{نیچے} \ur{لڑھکا۔}
\end{quote}

\begin{quote}
\emph{ek patthar ḍhalān se nīche luṛhkā.}
\end{quote}

\begin{quote}
A stone rolled down the slope.
\end{quote}

Past (it …): \ur{وہ} \ur{گیند} \ur{نیچے} \ur{لڑھکی}, \emph{voh gend nīche luṛhkī}.

\begin{cousinweb}
\textbf{\ur{لڑھکنا}} (\emph{luṛhaknā}) is inherited: CDIAL lists the Hindi cousin \emph{luṛhakṇā}, \textquotedblleft{}to roll, to slip\textquotedblright{}, under Sanskrit \emph{luṭhati}, \textquotedblleft{}rolls\textquotedblright{} (CDIAL 11079).
\end{cousinweb}

\begin{grammarlens}[title={What you've built}]
One more word for getting around.
\end{grammarlens}

\subsection*{Your turn}
\begin{itemize}
\raggedright
  \item \emph{Say it:} \emph{luṛhaknā}
  \item \emph{Say it:} \emph{luṛhaknā}, once more
  \item \emph{Say it:} \emph{ek patthar ḍhalān se nīche luṛhkā.} — A stone rolled down the slope.
  \item \emph{Recall:} say the Urdu for to show, then the Urdu for a slope, then say \emph{luṛhaknā} again
\end{itemize}

\begin{tcolorbox}[breakable,colback=teal!4,colframe=teal!35!black,title={Before you move on}]
What does \ur{لڑھکنا} mean? (\textquotedblleft{}To roll (down)\textquotedblright{}.) Say it once more.
\end{tcolorbox}
\section[\texorpdfstring{\ur{برآمدہ} (barāmadā)}{برآمدہ (barāmadā)}]{\ur{برآمدہ} (barāmadā) — a veranda}

\label{lesson:UR-C411-baramada}



\begin{quote}
Before the new one: say the Urdu for a slope, then the Urdu for to roll (down).
\end{quote}

\subsection*{You'll want to know: \ur{برآمدہ}}
\textbf{\ur{برآمدہ}} — \emph{barāmadā} — \textquotedblleft{}a veranda\textquotedblright{}.

\begin{quote}
\ur{ہم} \ur{برآمدے} \ur{میں} \ur{بیٹھ} \ur{کر} \ur{چائے} \ur{پیتے} \ur{ہیں۔}
\end{quote}

\begin{quote}
\emph{ham barāmade meṅ baiṭh kar chāe pīte haiṅ.}
\end{quote}

\begin{quote}
We sit on the veranda and drink tea.
\end{quote}

Masculine.

\begin{cousinweb}
\textbf{\ur{برآمدہ}} (\emph{barāmadā}) is usually traced to Persian \emph{barāmada}, \textquotedblleft{}come out, projecting\textquotedblright{}: the part of a house that juts out.
\end{cousinweb}

\begin{grammarlens}[title={What you've built}]
One more word for getting around.
\end{grammarlens}

\subsection*{Your turn}
\begin{itemize}
\raggedright
  \item \emph{Say it:} \emph{barāmadā}
  \item \emph{Say it:} \emph{barāmadā}, once more
  \item \emph{Say it:} \emph{ham barāmade meṅ baiṭh kar chāe pīte haiṅ.} — We sit on the veranda and drink tea.
  \item \emph{Recall:} say the Urdu for a slope, then the Urdu for to roll (down), then say \emph{barāmadā} again
\end{itemize}

\begin{tcolorbox}[breakable,colback=teal!4,colframe=teal!35!black,title={Before you move on}]
What does \ur{برآمدہ} mean? (\textquotedblleft{}A veranda\textquotedblright{}.) Say it once more.
\end{tcolorbox}
\section[\texorpdfstring{\ur{تہہ} \ur{خانہ} (tah khānā)}{تہہ خانہ (tah khānā)}]{\ur{تہہ} \ur{خانہ} (tah khānā) — a basement; a cellar}

\label{lesson:UR-C411-tah-khana}



\begin{quote}
Before the new one: say the Urdu for to roll (down), then the Urdu for a veranda.
\end{quote}

\subsection*{You'll want to know: \ur{تہہ} \ur{خانہ}}
\textbf{\ur{تہہ} \ur{خانہ}} — \emph{tah khānā} — \textquotedblleft{}a basement; a cellar\textquotedblright{}.

\begin{quote}
\ur{پرانی} \ur{عمارت} \ur{میں} \ur{ایک} \ur{تہہ} \ur{خانہ} \ur{ہے۔}
\end{quote}

\begin{quote}
\emph{purānī imārat meṅ ek tah khānā hai.}
\end{quote}

\begin{quote}
The old building has a basement.
\end{quote}

Masculine.

\begin{cousinweb}
\textbf{\ur{تہہ} \ur{خانہ}} (\emph{tah khānā}) is built from Persian \emph{tah}, \textquotedblleft{}the bottom\textquotedblright{}, and \emph{khāna}, \textquotedblleft{}a room\textquotedblright{}: it is usually traced to Persian \emph{tahkhāna}.
\end{cousinweb}

\begin{grammarlens}[title={What you've built}]
One more word for getting around.
\end{grammarlens}

\subsection*{Your turn}
\begin{itemize}
\raggedright
  \item \emph{Say it:} \emph{tah khānā}
  \item \emph{Say it:} \emph{tah khānā}, once more
  \item \emph{Say it:} \emph{purānī imārat meṅ ek tah khānā hai.} — The old building has a basement.
  \item \emph{Recall:} say the Urdu for to roll (down), then the Urdu for a veranda, then say \emph{tah khānā} again
\end{itemize}

\begin{tcolorbox}[breakable,colback=teal!4,colframe=teal!35!black,title={Before you move on}]
What does \ur{تہہ} \ur{خانہ} mean? (\textquotedblleft{}A basement; a cellar\textquotedblright{}.) Say it once more.
\end{tcolorbox}
\section[\texorpdfstring{\ur{بغل} (baghal)}{بغل (baghal)}]{\ur{بغل} (baghal) — beside; the side (of someone)}

\label{lesson:UR-C411-baghal}



\begin{quote}
Before the new one: say the Urdu for a veranda, then the Urdu for a basement.
\end{quote}

\subsection*{You'll want to know: \ur{بغل}}
\textbf{\ur{بغل}} — \emph{baghal} — \textquotedblleft{}beside; the side (of someone)\textquotedblright{}.

\begin{quote}
\ur{ہوٹل} \ur{کی} \ur{بغل} \ur{میں} \ur{ایک} \ur{دکان} \ur{ہے۔}
\end{quote}

\begin{quote}
\emph{hoṭal kī baghal meṅ ek dukān hai.}
\end{quote}

\begin{quote}
There's a shop next to the hotel.
\end{quote}

Feminine.

\begin{cousinweb}
\textbf{\ur{بغل}} (\emph{baghal}) is usually traced to Persian \emph{baghal}, \textquotedblleft{}the armpit, the side\textquotedblright{}; \textbf{\ur{کی} \ur{بغل} \ur{میں}} is \textquotedblleft{}right beside\textquotedblright{}.
\end{cousinweb}

\begin{grammarlens}[title={What you've built}]
One more word for getting around.
\end{grammarlens}

\subsection*{Your turn}
\begin{itemize}
\raggedright
  \item \emph{Say it:} \emph{baghal}
  \item \emph{Say it:} \emph{baghal}, once more
  \item \emph{Say it:} \emph{hoṭal kī baghal meṅ ek dukān hai.} — There's a shop next to the hotel.
  \item \emph{Recall:} say the Urdu for a veranda, then the Urdu for a basement, then say \emph{baghal} again
\end{itemize}

\begin{tcolorbox}[breakable,colback=teal!4,colframe=teal!35!black,title={Before you move on}]
What does \ur{بغل} mean? (\textquotedblleft{}Beside; the side (of someone)\textquotedblright{}.) Say it once more.
\end{tcolorbox}
